\section{Conclusion}
We have introduced LeanDojo: an open-source playground for learning-based theorem proving in Lean, consisting of toolkits, models, and benchmarks. It extracts data from Lean and enables the model to interact with Lean programmatically. We have developed {\name}, the first retrieval-augmented LLM for theorem proving. Limitations and future work are discussed in Appendix~\ref{appendix:limitations}.

We have released our code, data, models, and documentation to facilitate future research:

\begin{itemize}
    \item LeanDojo's codebase for data extraction and interaction with Lean: \url{https://github.com/lean-dojo/LeanDojo}
    \item LeanDojo's documentation: \url{https://leandojo.readthedocs.io}
    \item Datasets: (1) LeanDojo Benchmark: \url{https://doi.org/10.5281/zenodo.8016385} with DOI \texttt{10.5281/zenodo.8016385}. (2) LeanDojo Benchmark 4 (Appendix~\ref{appendix:lean4}): \url{https://doi.org/10.5281/zenodo.8040109} with DOI \texttt{10.5281/zenodo.8040109}.
    \item ReProver's code and models: \url{https://github.com/lean-dojo/ReProver}
    \item ChatGPT plugin (Appendix~\ref{appendix:chatgpt}): \url{https://github.com/lean-dojo/LeanDojoChatGPT}
    \item LeanDojo Website: \weburl
\end{itemize}
